\section{Qualitätssicherung der Anforderungen}

\subsection{Kriterien: Vollständigkeit, Eindeutigkeit, Konsistenz}

Alle funktionalen und nicht-funktionalen Anforderungen wurden anhand folgender Kriterien geprüft:

\begin{itemize}
    \item \textbf{Vollständigkeit}\\
    Es wurde sichergestellt, dass sämtliche sicherheitsrelevanten Funktionen (\ac{Hardstop}, \ac{Softstop}, \ac{E-Stop}, Quittierung, Resume, Kriechmodus, Warnsignale, Logging) vollständig beschrieben und durch passende Anforderungen abgedeckt sind.

    \item \textbf{Eindeutigkeit}\\
    Jede Anforderung wurde so formuliert, dass nur eine mögliche Interpretation zulässig ist. Unklare Begriffe wie „schnell“, „zuverlässig“, „rechtzeitig“ wurden vermieden oder messbar definiert.

    \item \textbf{Nachvollziehbarkeit}\\
    Anforderungen wurden so strukturiert, dass sie eindeutig auf Use-Cases, UML-Diagramme und Systemfunktionen zurückgeführt werden können. Jede wesentliche Systemfunktion findet sich in mindestens einem Modell wieder.

    \item \textbf{Konsistenz}\\
    Die Anforderungen wurden auf Widersprüche geprüft. Prioritätsregeln (\ac{Hardstop} > \ac{E-Stop} > \ac{Softstop}) wurden über alle Modelle hinweg einheitlich formuliert. Mehrfach vorkommende Anforderungen wurden zusammengeführt oder harmonisiert.

    \item \textbf{Verifizierbarkeit}\\
    Alle Anforderungen wurden so definiert, dass sie durch Tests, Simulationen oder statische Analyse überprüfbar sind. Insbesondere sicherheitsrelevante Anforderungen sind mit messbaren Bedingungen versehen.
\end{itemize}

\subsection{Vorgehen zur Qualitätssicherung}

Zur Sicherstellung der Qualität wurden folgende Schritte durchgeführt:

\begin{itemize}
    \item \textbf{Iterative Review-Zyklen}\\
    Anforderungen, Diagramme und Systembeschreibungen wurden mehrfach überarbeitet. Jede neue Modellierung (Use Case, Sequenz-, Zustands- oder Aktivitätsdiagramm) diente als Spiegel zur Überprüfung der Anforderungen.

    \item \textbf{Abgleich mit Use-Case-Modell}\\
    Für jeden Use-Case wurde geprüft, ob die entsprechenden funktionalen Anforderungen vollständig vorhanden sind.  
    Beispiele:
    \begin{itemize}
        \item „\ac{Hardstop}“ deckt FR-Hardstop ab
        \item „\ac{Softstop}“ deckt FR-Softstop und FR-Resume ab
        \item „Hindernis erkennen“ deckt FR-E-Stop ab
    \end{itemize}

    \item \textbf{Abgleich mit Zustandsdiagrammen}\\
    Die Zustandsmodelle dienten insbesondere dazu, die Vollständigkeit der Systemreaktionen zu prüfen 
    (z. B. Verhalten im \ac{Hardstop}, Übergänge im \ac{Softstop} und \ac{E-Stop}, Wiederaufnahmebedingungen).

    \item \textbf{Abgleich mit Sequenzdiagrammen}\\
    Die Sequenzdiagramme zeigen die tatsächlichen Nachrichtenflüsse und Systeminteraktionen. Dadurch konnten unvollständige oder unpräzise Anforderungen identifiziert und angepasst werden.

    \item \textbf{Validierung durch Architekturabgleich}\\
    Die Anforderungen wurden mit der Systemarchitektur (ROS-Nodes, CAN-Brücke, Arduino-Mikrocontroller, Motorcontroller, Sensorik) abgeglichen.  
    Ziel war: jede Anforderung muss technisch realisierbar und logisch zum Systemaufbau passend sein.
\end{itemize}

\subsection{Verifikation der Anforderungen}

Da im Rahmen dieser Projektarbeit ausschließlich eine Software Requirements Specification erstellt wurde, erfolgt die Verifikation der Anforderungen rein dokumentarisch und modellbasiert. Es wurden keine praktischen Tests oder Implementierungen durchgeführt.

Die Verifikation basiert auf drei Maßnahmen:

\begin{enumerate}
    \item \textbf{Modellbasierte Prüfung}\\
    Die UML-Diagramme (Use Case, Aktivität, Zustand, Sequenz) wurden genutzt, um die Anforderungen auf logische Konsistenz, Vollständigkeit und korrektes Systemverhalten zu überprüfen.

    \item \textbf{Dokumentarische Konsistenzprüfung}\\
    Die Anforderungen wurden iterativ überarbeitet und auf Eindeutigkeit, Widerspruchsfreiheit und klare Zuordnung zu Systemfunktionen geprüft.

    \item \textbf{Abgleich mit der Systemarchitektur}\\
    Die Anforderungen wurden mit der modellierten Architektur verglichen, um sicherzustellen, dass alle Funktionen technisch sinnvoll abbildbar sind und alle Schnittstellen konsistent beschrieben werden.
\end{enumerate}

Diese Maßnahmen gewährleisten, dass die spezifizierten Anforderungen nachvollziehbar, konsistent und verifizierbar formuliert sind.

\subsection{Nachvollziehbarkeit (Traceability)}

Zur Sicherstellung der Nachvollziehbarkeit besteht eine direkte Beziehung zwischen:

\begin{itemize}
    \item funktionalen Anforderungen (FR-xx)
    \item Use Cases
    \item Zustandsdiagrammen
    \item Sequenzdiagrammen
    \item Architekturkomponenten
\end{itemize}

Jede sicherheitsrelevante Anforderung taucht in mindestens einem UML-Diagramm wieder auf.

Beispiele:

\begin{itemize}
    \item FR-HS-xx → Use Case „\ac{Hardstop} betätigen“, Sequenzdiagramm \ac{Hardstop}, Zustandsdiagramm
    \item FR-SS-xx → Use Case „\ac{Softstop} betätigen“, \ac{Softstop}-Sequenzdiagramm, Zustandsdiagramm
    \item FR-ES-xx → Use Case „Hindernis erkennen (\ac{E-Stop})“, \ac{E-Stop}-Sequenzdiagramm, Aktivitätsdiagramm
    \item FR-RM-xx → Substatemaschinen \ac{Softstop}/\ac{E-Stop} + Resume-Sequenzdiagramm
\end{itemize}

Damit ist gewährleistet, dass jede Anforderung eindeutig auf modellierte Systemfunktionen rückführbar ist.
